%%%%%%%%%%%%%%%%%%%%%%%%%%%%%%%%%%%%%%%%%
% Arsclassica Article
% LaTeX Template
% Version 1.1 (1/8/17)
%
% This template has been downloaded from:
% http://www.LaTeXTemplates.com
%
% Original author:
% Lorenzo Pantieri (http://www.lorenzopantieri.net) with extensive modifications by:
% Vel (vel@latextemplates.com)
%
% License:
% CC BY-NC-SA 3.0 (http://creativecommons.org/licenses/by-nc-sa/3.0/)
%
%%%%%%%%%%%%%%%%%%%%%%%%%%%%%%%%%%%%%%%%%

%----------------------------------------------------------------------------------------
%	PACKAGES AND OTHER DOCUMENT CONFIGURATIONS
%----------------------------------------------------------------------------------------

\documentclass[
10pt, % Main document font size
a4paper, % Paper type, use 'letterpaper' for US Letter paper
oneside, % One page layout (no page indentation)
%twoside, % Two page layout (page indentation for binding and different headers)
headinclude,footinclude, % Extra spacing for the header and footer
BCOR5mm, % Binding correction
]{scrartcl}

\input{structure.tex} % Include the structure.tex file which specified the document structure and layout

\hyphenation{Fortran hy-phen-ation} % Specify custom hyphenation points in words with dashes where you would like hyphenation to occur, or alternatively, don't put any dashes in a word to stop hyphenation altogether
\usepackage[english]{babel}

%\usepackage{background}

\usepackage{graphicx}

\usepackage[os=win]{menukeys}

\usepackage{float}

\usepackage[dvipsnames,svgnames]{xcolor}

\usepackage[most]{tcolorbox}

\usepackage{listings}

\usepackage{caption}

\usepackage{enumitem}

%\usepackage{fontawesome}

\usepackage{fontawesome7}

\usepackage{dirtree}

\usepackage{textcomp}

\usepackage{hyperref}

\usepackage{svg}

%----------------------------------------------------------------------------------------
%	SETTINGS
%----------------------------------------------------------------------------------------



%\lstdefinestyle{DOS}
%{
%	backgroundcolor=\color{black},
%	basicstyle=\scriptsize\color{white}\ttfamily
%}

\lstdefinelanguage{DOS}
{
	backgroundcolor=\color{black},
	basicstyle=\scriptsize\color{white}\ttfamily,
	stringstyle=\color{white},
	breaklines=true,
}

\lstdefinelanguage{POWERSHELL}
{
	backgroundcolor=\color{DarkBlue},
	basicstyle=\scriptsize\color{white}\ttfamily,
	stringstyle=\color{white},
	breaklines=true,
    emph={New,Service},
	emphstyle={\color{yellow}},
}

\lstdefinelanguage{XML}
{
	basicstyle=\ttfamily\footnotesize,
	backgroundcolor = \color{LightGrey!10},
	basicstyle=\ttfamily\footnotesize,
	morestring=[b]",
	moredelim=[s][\bfseries\color{Maroon}]{<}{\ },
	moredelim=[s][\bfseries\color{Maroon}]{</}{>},
	moredelim=[l][\bfseries\color{Maroon}]{/>},
	moredelim=[l][\bfseries\color{Maroon}]{>},
	morecomment=[s]{<?}{?>},
	morecomment=[s]{<!--}{-->},
	commentstyle=\color{DarkOliveGreen},
	stringstyle=\color{blue},
	showstringspaces=false,
	identifierstyle=\color{red},
	breaklines=true,
	numbers=left
}

\definecolor{codegreen}{rgb}{0,0.6,0}
\definecolor{codegray}{rgb}{0.5,0.5,0.5}
\definecolor{codepurple}{HTML}{C42043}
\definecolor{backcolour}{HTML}{F2F2F2}
\definecolor{bookColor}{cmyk}{0,0,0,0.90}  
\color{bookColor}

\lstdefinelanguage{SQL}	
{
	basicstyle=\footnotesize\ttfamily,
    backgroundcolor=\color{backcolour},   
	commentstyle=\color{green},
	keywordstyle=\color{Maroon},
	numberstyle=\numberstyle,
	stringstyle=\color{codepurple},
	breakatwhitespace=false,
	breaklines=true,
	captionpos=b,
	keepspaces=true,
	numbers=left,
	numbersep=10pt,
	showspaces=false,
	showstringspaces=false,
	showtabs=false,
}

\newcommand\numberstyle[1]{%
	\footnotesize
	\color{codegray}%
	\ttfamily
	\ifnum#1<10 0\fi#1 |%
}

%\renewcommand{\contentsname}{My Custom Table of Contents}

\setlength{\parindent}{0pt}

%----------------------------------------------------------------------------------------
%	RENEWS
%----------------------------------------------------------------------------------------

\setlocalecaption{english}{contents}{TABLA DE CONTENIDOS}
\setlocalecaption{english}{table}{Tabla}
\setlocalecaption{english}{figure}{Figura}
\addto\captionsenglish{%
	\renewcommand\listfigurename{Listado de imágenes}}


\NewTotalTCBox{\commandbox}{ s v }
{verbatim,colupper=white,colback=black!75!white,colframe=black}
{\IfBooleanT{#1}{\textcolor{red}{\ttfamily\bfseries > }}%
	\lstinline[language=command.com,keywordstyle=\color{blue!35!white}\bfseries]^#2^}
	
\newtcolorbox{mybox}[1]{colback=yellow!5!white,
	colframe=yellow!65!black,fonttitle=\bfseries,
	title={#1},sharp corners}
	
	
\tcbset {
	base/.style={
		arc=0mm, 
		bottomtitle=0.5mm,
		boxrule=0mm,
		colbacktitle=black!10!white, 
		coltitle=black, 
		fonttitle=\bfseries, 
		left=2.5mm,
		leftrule=1mm,
		right=3.5mm,
		title={#1},
		toptitle=0.75mm, 
	}
}

\definecolor{brandblue}{rgb}{0.34, 0.7, 1}
\newtcolorbox{InfoBox}[1]{
	colframe=brandblue, 
	base={#1}
}

\newtcolorbox{WarningBox}[1]{
	%colback=yellow!5!white,
	colframe=yellow!60!white,
	base={#1}
}

\newtcolorbox{ErrorBox}[1]{
	%colback=yellow!5!white,
	colframe=red!60!white,
	base={#1}
}

\newtcolorbox{subbox}[1]{
	colframe=yellow!30!white,
	base={#1}
}	

\newtcbox{\xmybox}[1][red]{on line,
	arc=7pt,colback=#1!90!white,colframe=#1!50!black,
	before upper={\rule[-3pt]{0pt}{10pt}},boxrule=1pt,
	boxsep=0pt,left=6pt,right=6pt,top=2pt,bottom=2pt}


\renewcommand*\DTstylecomment{\rmfamily\color{JungleGreen}\textsc}
\renewcommand*\DTstyle{\ttfamily\textcolor{Maroon}}
%----------------------------------------------------------------------------------------
%	TITLE AND AUTHOR(S)
%----------------------------------------------------------------------------------------

\title{\normalfont\spacedallcaps{Phrasal Verb - Show Up}} % The article title

%\subtitle{Subtitle} % Uncomment to display a subtitle

\author{Rodrigo Castillejos} % The article author(s) - author affiliations need to be specified in the AUTHOR AFFILIATIONS block

\date{} % An optional date to appear under the author(s)

%----------------------------------------------------------------------------------------

\begin{document}

%----------------------------------------------------------------------------------------
%	HEADERS
%----------------------------------------------------------------------------------------

\renewcommand{\sectionmark}[1]{\markright{\spacedlowsmallcaps{#1}}} % The header for all pages (oneside) or for even pages (twoside)
%\renewcommand{\subsectionmark}[1]{\markright{\thesubsection~#1}} % Uncomment when using the twoside option - this modifies the header on odd pages
\lehead{\mbox{\llap{\small\thepage\kern1em\color{halfgray} \vline}\color{halfgray}\hspace{0.5em}\rightmark\hfil}} % The header style

\pagestyle{scrheadings} % Enable the headers specified in this block

%----------------------------------------------------------------------------------------
%	TABLE OF CONTENTS & LISTS OF FIGURES AND TABLES
%----------------------------------------------------------------------------------------

\maketitle % Print the title/author/date block

\setcounter{tocdepth}{3} % Set the depth of the table of contents to show sections and subsections only

\tableofcontents % Print the table of contents

%\newpage
%\listoffigures % Print the list of figures

%\listoftables % Print the list of tables

%----------------------------------------------------------------------------------------
%	ABSTRACT
%----------------------------------------------------------------------------------------

\newpage
%\section*{Resumen} % This section will not appear in the table of contents due to the star (\section*)

%----------------------------------------------------------------------------------------
%	AUTHOR AFFILIATIONS
%----------------------------------------------------------------------------------------

%\let\thefootnote\relax\footnotetext{* \textit{Department of Biology, University of Examples, London, United Kingdom}}

%\let\thefootnote\relax\footnotetext{\textsuperscript{1} \textit{Department of Chemistry, University of Examples, London, United Kingdom}}

%----------------------------------------------------------------------------------------

\newpage % Start the article content on the second page, remove this if you have a longer abstract that goes onto the second page

%----------------------------------------------------------------------------------------
%	DESINSTALACIÓN DE COMPONENTES
%----------------------------------------------------------------------------------------

\section{SHOW UP}
El \textit{phrasal verb} \textbf{show up} tiene un matiz diferente al verbo \textit{arrive} (llegar). mientras que \textit{arrive} es simplemente trasladarse del punto A al punto B, \textbf{show up} significa "aparecer", "presentarse" o "hacer acto de presencia" en un lugar.

\section{Conjugación del verbo (show)}

\begin{description}
	\item[Base/Presente] Show up (I show up)
	\item[Tercera persona presente] Shows up (He/She/It shows up)
	\item[Pasado simple] Showed up (I showep up)
	\item[Gerundio/Continuo] Showing up (I am showing up)
	\item[Participio] Shown up (I have shown up)
\end{description}

\begin{WarningBox}{\faIcon{triangle-exclamation} \, Advertencia} 
	\textbf{La regla del infinitivo:} Recuerda que cuando usas la palabra \textbf{"to"} como propósito de conexión de verbos (Ejemplo: I want to / I expect to), el verbo debe ir en su forma base, incluso si la frase habla en pasado. 
\end{WarningBox}

\begin{itemize}
	\item \textcolor{red}{\faIcon{xmark}} \textit{I expected him to showed up}
	\item \textcolor{green}{\faIcon{circle-check}} \textit{I expected him to \textbf{show up}}
\end{itemize}

\newpage
\section{Reto de traducción}

\begin{enumerate}
	\item \textbf{Él no apareció en la reunión esta mañana}
	\begin{itemize}
		\item \textit{He didn't show up to the meeting this morning.}
	\end{itemize}
	\item \textbf{Siempre apareces tarde cuando necesitamos ayuda.}
	\begin{itemize}
		\item \textit{You always show up late when we need help.}
	\end{itemize}
	\item \textbf{Estuvimos esperando por dos horas, pero el técnico nunca llegó.}
	\begin{itemize}
		\item \textit{We were waiting for two hours, but the technicion never showed up.}
	\end{itemize}
	\item \textbf{Asegúrate de presentarte a la entrevista con un traje.}
	\begin{itemize}
		\item \textit{Make sure to show up to the interview wearing a suit.}
	\end{itemize}
	\item \textbf{No esperaba que ella apareciera en la fiesta}
	\begin{itemize}
		\item \textit{I didn't expect her to show up to the party.}
	\end{itemize}
	\item \textbf{Siempre me pongo nervioso cuando el jefe aparece sin avisar.}
	\begin{itemize}
		\item \textit{I always get nervous when the boss shows up unannounced.}
	\end{itemize}
	\item \textbf{Si no apareces a tiempo, perderemos nuestra reservación.}
	\begin{itemize}
		\item \textit{If you don't show up on time, we will lose our reservation.}
	\end{itemize}
	\item \textbf{¿A qué hora apareció ella en la oficina hoy?}
	\begin{itemize}
		\item \textit{What time did she show up at the office today?}
	\end{itemize}
	\item \textbf{Es muy grosero invitar a alguien y luego no aparecer.}
	\begin{itemize}
		\item \textit{It's very rude to invite someone and then not show up.}
	\end{itemize}
	\item \textbf{Me sorprendió que tanta gente apareciera en el evento.}
	\begin{itemize}
		\item \textit{I was surprised that so many people showed up at the event.}
	\end{itemize}
\end{enumerate}
%----------------------------------------------------------------------------------------
%	BIBLIOGRAPHY
%----------------------------------------------------------------------------------------

%\renewcommand{\refname}{\spacedlowsmallcaps{References}} % For modifying the bibliography heading

%\bibliographystyle{unsrt}

%\bibliography{sample.bib} % The file containing the bibliography

%----------------------------------------------------------------------------------------

\end{document}